\documentclass[11pt]{article}
\usepackage[a4paper,margin=24mm]{geometry}
\usepackage{amsmath,amssymb}

\newcommand{\E}{\mathbb E}
\newcommand{\argmin}{\operatorname*{arg\,min}}
\setlength{\parindent}{0pt}
\setlength{\parskip}{5pt}
\setlength{\abovedisplayskip}{8pt}
\setlength{\belowdisplayskip}{8pt}

\begin{document}
\begin{center}
  {\Large\bfseries The first IRMF component}\par\vspace{4pt}
  {\large Notation and main results for smooth absolute loss}
\end{center}

At each location, the first IRMF component fits one value to nearby noisy
observations.  The theory compares that fitted value with an ideal population
fit and with the underlying clean signal.  This summary omits the proofs.

{\small
\renewcommand{\arraystretch}{1.12}
\begin{tabular}{@{}p{27mm}p{\dimexpr\linewidth-27mm-2\tabcolsep\relax}@{}}
\textbf{Symbol} & \textbf{Meaning} \\
\hline
\(u,\ t\) & \(u\) is an observation location; \(t\) is the location where we
compute the component.  Locations lie on the time axis \([0,1]\). \\
\(n,\ m\) & \(n\) is the number of observations; \(m\) is the number of
locations where we compute the component.  Often \(m=n\). \\
\(X(\cdot),X_u,X_t\) & \(X\) is the unknown clean signal.  \(X_u=X(u)\) and
\(X_t=X(t)\) are its values at the observation and fitting locations. \\
\(Y_u\) & The observed noisy value at location \(u\). \\
\(\varepsilon_u,\sigma\) & \(\varepsilon_u=Y_u-X_u\) is observation noise;
\(\sigma\) is its standard deviation, in signal-amplitude units. \\
\(h>0\) & Window radius, in time units.  Only observations with
\(|u-t|\leq h\) contribute to the fit at \(t\). \\
\(H>0\) & Smoothing scale of the loss, in signal-amplitude units.  It
controls how closely the loss follows the absolute value. \\
\(a_{u,t}\) & The weight of observation \(u\) in the fit at \(t\).
Weights are nonnegative, fixed before seeing the noise, and sum to one. \\
\(\widetilde S_t\) & The component fitted to the observed data.  It changes
when the noise realization changes. \\
\(S_t\) & The population component: the minimizer of the loss averaged over
repeated noise realizations, for fixed signal, weights, loss scale, and noise law. \\
\end{tabular}}

The observation model is
\(Y_u=X_u+\varepsilon_u\), with independent Gaussian errors
\(\varepsilon_u\sim N(0,\sigma^2)\), meaning mean zero and variance
\(\sigma^2\).  For example, \(t=0.40\) and \(h=0.05\) use observations in
\([0.35,0.45]\).

The loss \(\rho_H\) penalizes an amplitude residual \(r\).
Its derivative \(\psi_H\), called the score, measures an observation's
contribution to the fitting equation:
\[
  \rho_H(r)=\E_Z|r+HZ|,\qquad \psi_H(r)=\rho_H'(r),\qquad |\psi_H(r)|\leq1.
\]
Here \(Z\sim N(0,1)\) is an auxiliary standard Gaussian variable, independent
of the observation noise.  \(\E\) means expectation, or averaging over the
specified random variable.  Primes mean derivatives with respect to the
residual.  The bounded score limits each observation's contribution.

In the definitions below, \(z\) is a candidate fitted amplitude, so
\(Y_u-z\) is a residual.  The notation \(\argmin_z\) means the value of
\(z\) that minimizes the displayed sum, and \(\sum_u\) sums over observations:
\begin{align*}
  \widetilde S_t&=\argmin_z\sum_u a_{u,t}\rho_H(Y_u-z),
  & S_t&=\argmin_z\sum_u a_{u,t}\E\rho_H(Y_u-z).
\end{align*}
The expectation defining \(S_t\) is over observation noise.  For Gaussian
noise, let \(B\) denote the effective loss scale after that averaging:
\[
  B=\sqrt{H^2+\sigma^2},\qquad
  S_t=\argmin_z\sum_u a_{u,t}\rho_B(X_u-z).
\]
Here \(\rho_B\) is the same loss with \(H\) replaced by \(B\).  Thus
\(X_t\) is the clean signal value, \(S_t\) is an ideal local fit, and
\(\widetilde S_t\) is the fit computed from noisy data.  They generally differ.

\newpage

Three more quantities describe the estimation error at \(t\):

{\small
\renewcommand{\arraystretch}{1.1}
\begin{tabular}{@{}p{27mm}p{\dimexpr\linewidth-27mm-2\tabcolsep\relax}@{}}
\(F_t>0\) & Population curvature: how strongly the expected loss bends upward
near its minimum \(S_t\).  It is a scalar measuring local stability. \\
\(\nabla_t\) & Score fluctuation caused by the observed noise, after
subtracting its expectation.  Despite the symbol, it is a scalar here. \\
\(R_t\) & The remainder left after subtracting the leading random error
\(\nabla_t/F_t\) from \(\widetilde S_t-S_t\). \\
\end{tabular}}

Their definitions and the main expansion are
\begin{align*}
  F_t&=\sum_u a_{u,t}\E\rho_H''(Y_u-S_t)>0,\\
  \nabla_t&=\sum_u a_{u,t}
    \bigl\{\psi_H(Y_u-S_t)-\E\psi_H(Y_u-S_t)\bigr\}.
\end{align*}
\[
  \boxed{\widetilde S_t-S_t=\frac{\nabla_t}{F_t}+R_t.}
\]
Proposition 4.1 bounds \(R_t\).  Lemma 4.2 controls the score and curvature
fluctuations used in that bound.  The corrected value
\(\widetilde S_t^{\mathrm{corr}}=\widetilde S_t-\nabla_t/F_t\) has error
\(R_t\) relative to \(S_t\).  The superscript ``corr'' labels this correction.
It is an \emph{oracle} correction because it uses the unknown \(S_t\) and
population quantities.

The effective sample size is \(n_{\mathrm{eff}}(t)=1/\sum_u a_{u,t}^2\).
For regular weights and enough grid points per window, it is of order
\(h/\Delta\), where \(\Delta\) is the grid spacing.  On \([0,1]\), the
usual grids have \(\Delta=1/n\) or \(1/(n-1)\), both giving order \(nh\).
Each individual weight is then at most a constant times \(1/(nh)\).

Choose a failure probability \(0<\beta<1\), and define the logarithmic factor
\(\ell=\log(4m/\beta)\), using the natural logarithm.  For a 95\% guarantee,
take \(\beta=0.05\).  Keep \(H\), \(\sigma\), and the weight profile fixed;
assume the population curvatures \(F_t\) stay above a common positive constant.
For sufficiently large \(nh/\ell\), with probability at least \(1-\beta\),
\[
  \boxed{\begin{aligned}
    \max_t|\widetilde S_t-S_t|
      &\lesssim\sqrt{\frac{\ell}{nh}},\\[3pt]
    \max_t|R_t|&\lesssim\frac{\ell}{nh}.
  \end{aligned}}
\]
Here \(\max_t\) means the largest error across the \(m\) fitting locations.
The symbol \(\lesssim\) means ``at most a constant times''; the constants
depend on \(H\), \(\sigma\), the weight profile, and the curvature lower
bound.  These are upper bounds on the uncorrected and oracle-corrected errors.

The population bias depends on the window radius through the population fit:
\(b_t(h)=S_t(h)-X_t\), with \(H\) held fixed.  If the signal has bounded
first and second derivatives and equal weights at equal distances on opposite
sides of \(t\), Proposition 4.5 gives \(\max_t|b_t(h)|\lesssim h^2\).
Thus \(h^2\) bounds the bias magnitude; it is not the bias itself.
Under these conditions, the same probability guarantee gives
\[
  \max_t|\widetilde S_t-X_t|
  \lesssim
  \underbrace{h^2}_{\text{bias bound}}
  +\underbrace{\sqrt{\ell/(nh)}}_{\text{random error}}
  +\underbrace{\ell/(nh)}_{\text{remainder}}.
\]
Smaller \(h\) tightens the bias bound but raises the random-error bound.
Balancing them with \(h\) proportional to \((\ell/n)^{1/5}\) gives a bound
of order \((\ell/n)^{2/5}\), as \(\ell/n\to0\).

The results concern the first component on a finite grid with enough local
data.  Boundary truncation can break symmetry.  Asymmetric contamination,
later recursive components, and choices of \(H\) that change with sample
size need additional checks.

{\small See \texttt{proposition\_4\_1\_lemma\_4\_2.tex} for the constants
and sample-size conditions, with proofs in Appendix A.  This summary uses
normalized weights; \(x=\log(1/\beta)\) and \(z_x=\ell\) in the full note.}

\end{document}
